%==============================================================================
% Section 2: Lorentzian EC+NY setup and Bianchi class-A backgrounds
%==============================================================================
\section{Lorentzian EC+NY setup and Bianchi class-A backgrounds}
\label{sec:setup}

\subsection{First-order action ingredients}
\label{sec:setup:action}

The basic variables are the tetrad one-form $e^{I}$ and the Lorentz connection
one-form $\omega^{IJ}$, both treated as independent (first-order / Palatini
type).  The action consists of the Einstein--Cartan term~\cite{ref8} and the
Nieh--Yan term~\cite{ref9}:
\begin{equation}
  S[e,\omega]
  =\frac{1}{2\kap^{2}}\int \epsilon_{IJKL}\, e^{I}\wedge e^{J}\wedge R^{KL}[\omega]
  +\frac{\thetaNY}{\kap^{2}}\int \left( T^{I}\wedge T_{I}
       - e_{I}\wedge e_{J}\wedge R^{IJ} \right),
\end{equation}
where $R^{IJ}[\omega]=\dd\omega^{IJ}+\omega^{I}{}_{K}\wedge\omega^{KJ}$ and
$T^{I}=\dd e^{I}+\omega^{I}{}_{J}\wedge e^{J}$.  The signature and the
Hodge/$\epsilon$ conventions are fixed in Appendix~\ref{app:conventions}.

\paragraph{Treatment of the NY term.}
We adopt the route that treats the NY term as an exact four-form (the NY-MAIN
route).  Within our NY-MAIN route and the inspected reduced/native-action gate
of the LZ-native $\AX/\VT/\MX$ torsion ansatz scalars, the NY term contributes
on the background only as an exact boundary term.  In the reduced system it
forms an exact auxiliary shell, and its promotion to fields $\eta(x),V(x)$
generates no quadratic kinetic/gradient term under the inspected native-action
criterion (Section~\ref{sec:torsion}).  The topological nature of the NY term in
the canonical analysis is known~\cite{ref33,ref34,ref35}.  This statement of
bulk non-involvement is restricted to the above scope and does not extend
automatically to a full torsion field theory or to a general NY action.

\subsection{Bianchi class-A backgrounds and the structure matrix}
\label{sec:setup:structure}

The spatial sections carry the left-invariant structure of a three-dimensional
unimodular Lie group $G$.  The left-invariant one-forms $\sigma^{i}$
($i=1,2,3$) satisfy
\begin{equation}
  \dd\sigma^{i}=-\tfrac12\,C^{i}{}_{jk}\,\sigma^{j}\wedge\sigma^{k}.
\end{equation}
The class-A (unimodular) condition is $C^{j}{}_{ij}=0$~\cite{ref15}, and then the
structure constants are encoded in a symmetric \emph{structure matrix}
$n^{ij}$:
\begin{equation}
  C^{i}{}_{jk}=\epsilon_{jkl}\,n^{li},
  \qquad
  n^{ij}=\tfrac12\,\epsilon^{ikl}\,C^{j}{}_{kl},
  \qquad n^{ij}=n^{ji}.
\end{equation}
By the Milnor frame~\cite{ref16}, $n^{ij}=\mathrm{diag}(n_{1},n_{2},n_{3})$, and
its rank/signature classifies the class-A type.  The four topologies treated
here are collected in Table~\ref{tab:bianchi}.  Bianchi models with compact
spatial sections are restricted to class A~\cite{ref19,ref20}.

\begin{table}[H]
\centering
\begingroup
\renewcommand{\arraystretch}{1.4}
\begin{tabular}{llcrrrr}
\toprule
Topology & Bianchi & $n$ (Milnor frame) & $\rank(n)$ & $\dzero=3-\rank(n)$ & $b_{1}$ & $R\,q^{2}$ \\
\midrule
$\Tthree$   & I    & $\mathrm{diag}(0,0,0)$  & 0 & 3 & 3 & $0$    \\
$\Nilthree$ & II   & $\mathrm{diag}(1,0,0)$  & 1 & 2 & 2 & $-1/2$ \\
$\Solthree$ & VI$_{0}$ & $\mathrm{diag}(1,-1,0)$ & 2 & 1 & 1 & $-2$   \\
$\Sthree$   & IX   & $\mathrm{diag}(1,1,1)$  & 3 & 0 & 0 & $6$    \\
\bottomrule
\end{tabular}
\endgroup
\caption{Four Bianchi class-A backgrounds.  Here $b_{1}$ is the first Betti
number of the corresponding compact quotient (standard values~\cite{ref18}), and
$R\,q^{2}$ is the dimensionless spatial scalar curvature at the isotropic point.}
\label{tab:bianchi}
\end{table}

The $\Sthree$ structure-constant normalisation is fixed to $C=\epsilon$ ($|C|=1$)
throughout (Appendix~\ref{app:conventions}, where the explicit conversion to the
normalisation of paper~VI~\cite{ref7} is also recorded).  The quantities
$\rank(n)$ and $\dzero$ are invariant under an overall structure-constant scaling
$C\to sC$, whereas the \emph{numerical} values such as $R\,q^{2}$ and the TT
eigenvalues depend on $s$ (Appendix~\ref{app:conventions:s3}).

The $\dzero$ in Table~\ref{tab:bianchi} is the algebraic prediction from the
rank defect of the structure matrix.  That this is realised as the kernel of the
constrained spectrum is the main result of this paper
(Section~\ref{sec:rankdefect}).  The relation to abelianisation,
$\dim(\mathfrak{g}/[\mathfrak{g},\mathfrak{g}])=\dim\ker n$, is also a known
algebraic fact and is used only as a geometric cross-check.

\subsection{Background ansatz and perturbation sectors}
\label{sec:setup:perturbation}

The background is taken as the homogeneous diagonal ansatz
\begin{equation}
  e^{0}=N(t)\,\dd t,\qquad
  e^{a}=q_{a}(t)\,\sigma^{a} \quad (a=1,2,3;\ \text{no sum}),
\end{equation}
and the connection $\omega^{IJ}$ has independent homogeneous components (the
explicit component list and the reduced action are in
Appendix~\ref{app:native}).  Relative to the background frame the perturbations
decompose into

\begin{itemize}
  \item a \textbf{scalar (diagonal/axis) sector}: diagonal metric perturbations
        $\beta_{a}$ and the corresponding connection perturbations,
  \item a \textbf{spin-1 (transverse vector) sector},
  \item a \textbf{spin-2 (TT) sector}: transverse-traceless perturbations
        $h^{\mathrm{TT}}_{ab}$.
\end{itemize}

The analysis of the homogeneous modes ($k=0$) is central to
Section~\ref{sec:rankdefect}, while the inhomogeneous modes are treated by a
topology-dependent harmonic expansion ($\Tthree$: Fourier; $\Sthree$:
hyperspherical harmonics~\cite{ref28,ref29}; $\Nilthree$: toroidal $\cup$
Heisenberg sector; $\Solthree$: boost sector) in Section~\ref{sec:spectral} and
Appendix~\ref{app:spectral}.

\subsection{Torsion ansatz scalars}
\label{sec:setup:torsion}

On the homogeneous background, the torsion within the ansatz used here organises
into the following scalar modes:

\begin{itemize}
  \item \textbf{AX (axial)}: an axial scalar $\eta(t)$ giving
        $T^{a}{}_{bc}\supset\eta\,\epsilon_{abc}$;
  \item \textbf{VT (vector/trace)}: a scalar $V(t)$ giving the trace part;
  \item \textbf{MX (mixed)}: the mixed branch of the two.
\end{itemize}

These are the same family as the branch classification (EH/AX/VT/MX) of
papers~V/VI~\cite{ref6,ref7}.  Section~\ref{sec:torsion} classifies the
derivative couplings obtained when these are promoted to fields
$\eta(x),V(x)$.  All torsion claims of this paper lie \emph{within the
$\AX/\VT/\MX$ ansatz scalars}; no generalisation to a full 24-component torsion
field theory is claimed.

\subsection{Notation summary}
\label{sec:setup:notation}

In what follows, $q^{2}$ denotes the square of the isotropic scale factor (the
isotropic point $q_{1}=q_{2}=q_{3}=q$), $\kap^{2}=8\pi G$, and $\vol$ denotes
the comoving volume of the spatial cell.  The notation of the constraint
analysis (primary/secondary constraints, first/second class, Dirac bracket)
follows the standard definitions~\cite{ref10,ref11,ref12}.
